\subsection{Скатывания с горок}
%%%%%%%%%%%%%%%%%%%%%%%%%%%%%%%%%%%%%%%%%%%%%%%%%%%%
%%%%%%%%%%%%%%%%%%%%%%%%%%%%%%%%%%%%%%%%%%%%%%%%%%%%
\z[{\begin{wrapfigure}[7]{r}{0.2\textwidth}%
		\vspace{-12pt}%
		\begin{flushright}%
 			\includegraphics[width=0.2\textwidth]{7.3.1}%
		\end{flushright}%
	\end{wrapfigure}%
	}]%
	{Система из грузов $m$ и $M$ и связывающей их лёгкой нерастяжимой нити в начальный момент покоится в вертикальной плоскости, проходящей через центр закреплённой сферы. Груз $m$ находится на вершине сферы (см. рисунок). В ходе возникшего движения груз $m$ отрывается от поверхности сферы, пройдя по ней дугу $30\degree$. Найдите массу $M$, если $m = 100\,\text{г}$. Размеры груза $m$ ничтожно малы по сравнению с радиусом сферы. Трением пренебречь.}
%
{$M = m \cdot \frac{3\cos\alpha - 2}{2\alpha - \cos\alpha}, \quad \text{где } \alpha = \frac{\pi}{6}, \quad M \approx 329\,\text{г}$}
%%%%%%%%%%%%%%%%%%%%%%%%%%%%%%%%%%%%%%%%%%%%%%%%%%%% 
 \z{Небольшой шарик соскальзывает без начальной скорости и без трения с верхней точки сферы, закреплённой на горизонтальной поверхности стола. Под каким углом $\theta$ к поверхности стола шарик ударится о стол?}
%
{$\alpha = \arccos\left(\frac{\sqrt{6}}{9}\right)$}
%%%%%%%%%%%%%%%%%%%%%%%%%%%%%%%%%%%%%%%%%%%%%%%%%%%% 
\z[{\begin{wrapfigure}[6]{r}{0.35\textwidth}%
		\vspace{-12pt}%
		\begin{flushright}%
 			\includegraphics[width=0.35\textwidth]{7.3.3}%
		\end{flushright}%
	\end{wrapfigure}%
	}]%
	{Полушар радиусом $R$ покоится на горизонтальной поверхности стола. В точку $A$ на полушаре помещают небольшую по сравнению с размерами полушара шайбу массой $m$ и отпускают (см. рисунок). Шайба скользит без трения и оказывается в точке $B$, а полушар при этом остаётся неподвижным. Радиусы $OA$ и $OB$ составляют с вертикалью углы $\alpha$ и $\beta$, такие, что $\cos\alpha = \frac{5}{6}$, $\cos\beta = \frac{2}{3}$.
	\lpu Найти скорость шайбы в точке $B$. \lpuend
	\lpu Найти силу трения между полушаром и столом при прохождении шайбой точки $B$.  \lpuend
	}
%
{\puan  $v = \sqrt{\frac{gR}{3}};$~\puan $F_{\text{тр}} = \frac{mg\sqrt{5}}{9}$}
%%%%%%%%%%%%%%%%%%%%%%%%%%%%%%%%%%%%%%%%%%%%%%%%%%%%	
\z[{\begin{wrapfigure}[5]{r}{0.2\textwidth}%
		\vspace{-12pt}%
		\begin{flushright}%
 			\includegraphics[width=0.2\textwidth]{7.3.4}%
		\end{flushright}%
	\end{wrapfigure}%
	}]%
	{На гладком горизонтальном столе находится чаша массой $M$ с полусферической выемкой радиусом $R$ с гладкими стенками (смотри рисунок). На самый край выемки чаши поместили монету массой $m$, размеры которой значительно меньше размеров выемки. В начальный момент монета и чаша друг относительно друга не двигались. Монету и чашу одновременно отпустили. С каким ускорением движется монета, проходя самое нижнее положение?}
%
{$a = \frac{2(M + m)}{M}g$}
%%%%%%%%%%%%%%%%%%%%%%%%%%%%%%%%%%%%%%%%%%%%%%%%%%%%	
\z[{\begin{wrapfigure}[5]{r}{0.2\textwidth}%
		\vspace{-12pt}%
		\begin{flushright}%
 			\includegraphics[width=0.2\textwidth]{7.3.5}%
		\end{flushright}%
	\end{wrapfigure}%
	}]%
	{Тонкое проволочное кольцо массы $M$ стоит на горизонтальной плоскости (рис.). По кольцу могут скользить без трения две одинаковые бусинки массой $m$ каждая. В начальный момент времени бусинки находятся вблизи верхней точки кольца. Их одновременно отпускают, и они начинают двигаться симметрично. При каком отношении масс $n = m/M$ кольцо оторвётся от плоскости?}
%
{$n \geq \frac{3}{2}$}